% Part: first-order-logic
% Chapter: beyond
% Section: second-order-logic

\documentclass[../../../include/open-logic-section]{subfiles}

\begin{document}

\olfileid{fol}{byd}{sol}

\olsection{Logica van de tweede orde}

In de logica van de tweede orde mogen we over twee soorten dingen
kwantificeren. We kunnen niet alleen spreken over de individuen in
!!a{domain}, maar ook over relaties op dat !!{domain}. We beginnen met een
eerste-ordetaal $\Lang{L}$. Voor elk getal $k$ voegen we !!{variable}s $R$
toe. Zo'n variabele kan voor elke $k$-plaatsige relatie staan: een relatie met
zoveel plaatsen. Over zo'n variabele mogen we dan kwantificeren. Als $R$ voor
zo'n $k$-plaatsige relatie staat en $t_1$, \dots,~$t_k$ gewone termen van de
eerste orde zijn,
dan is $\Atom{R}{t_1,\dots,t_k}$ een atomaire !!{formula}. De overige
!!{formula}s maken we op dezelfde manier als in de eerste-ordelogica. We
voegen alleen regels toe voor kwantoren die over de nieuwe variabelen lopen.
!!{identity} hebben we rechtstreeks alleen voor termen van de eerste orde.
Zijn $R$ en $S$ allebei variabelen voor een $k$-plaatsige relatie, dan kunnen
we $\eq[R][S]$ wel gebruiken als afkorting voor
\[
\lforall[x_1 \dots][\lforall[x_k][(\Atom{R}{x_1, \dots, x_k} \liff
  \Atom{S}{x_1, \dots, x_k})]].
\]
Deze formule zegt dat de twee relaties voor precies dezelfde rijtjes van
objecten gelden.

De bewijsregels voor de tweede orde zijn in de eerste plaats de bekende
kwantorregels, maar dan ook voor de nieuwe tweede-ordevariabelen. We moeten wel
precies zeggen hoe die variabelen samengaan met de !!{predicate}s van
~$\Lang{L}$ en, algemener, met de !!{formula}s van~$\Lang{L}$. Om te beginnen tellen relatievariabelen als
termen. Daardoor mogen we bijvoorbeeld afleiden:
\[
!A(R) \Proves \lexists[R][!A(R)]
\]
Stel nu dat $\Lang{L}$ de taal van de rekenkunde is en dat $<$ daarin een
vast relatiesymbool is. Dan verwachten we ook dat deze afleiding geldig is:
\[
x < y \Proves \lexists[R][\Atom{R}{x,y}]
\]
Voor een gegeven !!{formula}~$!A$ verwachten we net zo goed:
\[
!A(x_1, \dots, x_k) \Proves \lexists[R][\Atom{R}{x_1,\dots,x_k}]
\]
In het algemeen zouden we de volgende afleidingen kunnen toelaten:
\[
\Subst{!A}{\lambd[\vec x][!B(\vec x)]}{R} \Proves
\lexists[R][!A]
\]
De uitdrukking $\Subst{!A}{\lambd[\vec x][!B(\vec x)]}{R}$ krijgen we als
we elke atomaire !!{formula} van de vorm $\Obj{R}{t_1,\dots,t_k}$ die in
~$!A$ voorkomt, vervangen door $!B(t_1, \dots, t_k)$. Deze regel
komt op hetzelfde neer als een {\em comprehensieschema}. Dat is een vast
patroon voor axioma's van de volgende vorm:
\[
\lexists[R][\lforall[x_1, \dots, x_k][(!A(x_1, \dots, x_k) \liff
\Atom{R}{x_1, \dots, x_k})]],
\]
We nemen zo'n axioma op voor elke !!{formula}~$!A$ van de tweede-ordetaal
waarin $R$ niet vrij voorkomt. (Opgave: laat zien dat dit schema tot een
tegenspraak leidt als $R$ wel in~$!A$ mag voorkomen.)

Met de ``axioma's van de tweede-ordelogica'' bedoelen logici meestal het
volgende. Ze nemen de eerste-ordelogica, voegen de kwantorregels voor de
tweede orde toe en nemen ook het comprehensieschema op. Dat is de minimale
uitbreiding. We kunnen ook zwakkere delen van deze axioma's en regels
bestuderen. Het volledige comprehensieschema is bijvoorbeeld
\emph{impredicatief}. Dat betekent hier het volgende. Een !!{formula} met
tweede-ordekwantoren mag een relatie $\Atom{R}{x_1, \dots, x_k}$
``definiëren''. Maar die kwantoren lopen al over de verzameling van alle
zulke relaties, dus ook over $R$ zelf!{} Rond 1900 reageerden veel mensen op
de paradox van Russell door juist zulke definities de schuld te geven. Bij de
opbouw van de grondslagen van de wiskunde probeerden ze deze definities daarom
te vermijden. Verbieden we tweede-ordekwantoren in de !!{formula}~$!A$, dan
krijgen we een {\em predicatieve} vorm van comprehensie. Die vorm is iets
zwakker.

Nu bekijken we de betekenis, dus de semantiek. Een tweede-orde!!{structure}
bestaat eerst uit een gewone eerste-orde!!{structure} voor de taal. Daar komt
een verzameling relaties op het !!{domain} bij. De tweede-ordekwantoren lopen
over die relaties. Preciezer gezegd hebben we voor elke $k$ een verzameling
$k$-plaatsige relaties. Staat comprehensie in het !!{derivation} systeem, dan
moet dit tweede deel genoeg relaties bevatten om de comprehensieaxioma's waar
te maken. Anders is het !!{derivation} systeem niet correct. We kunnen daar
gemakkelijk voor zorgen door in het tweede deel \emph{alle} relaties op het
eerste deel op te nemen. Zo'n !!a{structure} heet \emph{volledig}. In zekere
zin is dit de ``bedoelde !!{structure}'' voor de taal. Als we alleen volledige
!!{structure}s toelaten, gebruiken we de \emph{volledige semantiek} voor de
tweede-ordelogica. We hoeven dan alleen het eerste-ordedeel vast te leggen.
Welke relaties in het tweede-ordedeel zitten, ligt daarna al vast: het zijn ze
allemaal.

Er zijn dus verschillende dingen die iemand met ``tweede-ordelogica'' kan
bedoelen. Voor het !!{derivation} systeem zijn er twee mogelijkheden:
\begin{enumerate}
\item het ``minimale'' systeem voor de tweede orde, met een aantal
  comprehensieaxioma's;
\item het ``standaardsysteem'' voor de tweede orde, met volledige
  comprehensie.
\end{enumerate}
Ook voor de semantiek zijn er twee mogelijkheden:
\begin{enumerate}
\item de ``zwakke'' semantiek: de !!{structure} heeft een eerste-ordedeel en
  daarnaast een tweede-ordedeel dat groot genoeg is voor de gekozen
  comprehensieaxioma's;
\item de ``standaardsemantiek'' voor de tweede orde: we laten alleen
  volledige !!{structure}s toe, dus structuren met alle relaties.
\end{enumerate}
Als logici niet zeggen welk systeem of welke semantiek ze bedoelen, bedoelen
ze meestal de tweede mogelijkheid in beide lijsten. De standaardsemantiek
heeft een groot voordeel. Daarmee kunnen we veel bekende wiskundige
structuren categorisch beschrijven: de beschrijving legt de structuur tot op
isomorfie vast. Het !!{derivation} systeem is bovendien vrij sterk en is
correct voor deze semantiek. Er staat een nadeel tegenover. Het systeem is
\emph{niet} volledig voor de semantiek. Er bestaat zelfs \emph{geen}
effectief gegeven !!{derivation} systeem dat volledig is voor de volledige
tweede-ordesemantiek. Voor de zwakkere semantiek is het systeem \emph{wel} volledig.
Als een zin niet bewijsbaar is, bestaat er dus \emph{een} !!{structure}
waarin die zin onwaar is. Dat hoeft alleen niet de volledige structuur te
zijn.

Met tweede-ordelogica kunnen we veel uitdrukken. Een eenplaatsige relatie
kunnen we zien als een deelverzameling van het !!{domain}. Daardoor kunnen we
over zulke verzamelingen kwantificeren. We kunnen de inductie voor de
natuurlijke getallen bijvoorbeeld in één axioma zetten:
\[
\lforall[R][((\Atom{R}{\Obj{0}} \land \lforall[x][(\Atom{R}{x} \lif
    \Atom{R}{x'})]) \lif \lforall[x][\Atom{R}{x}])].
\]
Neem voor de taal van de rekenkunde de symbolen $\Obj 0, \Obj \prime, +,
\times$ en $<$. Met de volgende axioma's leggen we hun gedrag vast:
\begin{enumerate}
\item $\lforall[x][\lnot x' = \Obj 0]$
\item $\lforall[x][\lforall[y][(s(x) = s(y) \lif x = y)]]$
\item $\lforall[x][(x + \Obj 0) = x]$
\item $\lforall[x][\lforall[y][(x + y') = (x + y)']]$
\item $\lforall[x][(x \times \Obj 0) = \Obj 0]$
\item $\lforall[x][\lforall[y][(x \times y') = ((x \times y) + x)]]$
\item $\lforall[x][\lforall[y][(x < y \liff \lexists[z][y = (x + z')])]]$
\end{enumerate}
Gebruiken we de volledige tweede-ordesemantiek, dan beschrijven deze axioma's
samen met het inductieaxioma precies de !!{structure}~$\Struct{N}$, het
standaardmodel van de rekenkunde. ``Precies'' betekent hier: op isomorfie na.
We kunnen dat als volgt laten zien. Neem een willekeurige
!!{structure}~$\Struct{M}$ waarin de axioma's waar zijn. We definiëren een
functie~$f$ van $\Nat$ naar het !!{domain} van $\Struct{M}$. Daarvoor
gebruiken we gewone recursie op~$\Nat$. We zetten
$f(0) = \Assign{\Obj 0}{M}$ en
$f(x+1) = \Assign{\prime}{M}(f(x))$. Gewone inductie op~$\Nat$ laat samen
met het feit dat axioma's (1) en~(2) in~$\Struct M$ gelden, zien dat $f$
!!{injective} is. Nu moeten we nog laten
zien dat $f$ !!{surjective} is. Laat~$P$ bestaan uit de elementen van
$\Domain{M}$ die als waarde van~$f$ voorkomen. Omdat $\Struct M$ volledig is,
hoort $P$ bij het tweede-orde!!{domain}. Door de manier waarop we $f$ hebben
gemaakt, ligt $\Assign{\Obj 0}{M}$ in~$P$ en blijft een element in~$P$ als we
er $\Assign{\prime}{M}$ op toepassen. Het inductieaxioma geldt in $\Struct M$
en dus ook voor~$P$. Daarom is $P$~het hele eerste-orde!!{domain} van
$\Struct M$. De functie $f$ is dus !!a{bijection}. We kunnen net zo aantonen
dat $f$ een homomorfisme is, met herhaalde gewone inductie op $\Nat$.

In de verzamelingenleer is een functie een bijzonder soort relatie. Een
eenplaatsige functie~$f$ kunnen we bijvoorbeeld weergeven als een
tweeplaatsige relatie~$R$ waarvoor
$\lforall[x][\lexists![y][R(x,y)]]$ geldt. We mogen daardoor ook over
functies kwantificeren. Met de volledige semantiek kunnen we nu de klasse van
oneindige !!{structure}s $\Struct M$ definiëren. Daarbij moet een injectieve
functie bestaan van het !!{domain} van $\Struct M$ naar een echte
deelverzameling van datzelfde !!{domain}:
\[
\lexists[f][(\lforall[x][\lforall[y][(\eq[f(x)][f(y)] \lif \eq[x][y])]]
  \land \lexists[y][\lforall[x][\eq/[f(x)][y]]])].
\]
De klasse van eindige !!{structure}s wordt dan juist gedefinieerd door de
ontkenning van deze zin.

We kunnen ook welordeningen definiëren. Daarvoor voegen we aan de definitie
van een lineaire ordening de volgende voorwaarde toe:
\[
\lforall[P][(\lexists[x][\Atom{P}{x}] \lif \lexists[x][(\Atom{P}{x}
    \land \lforall[y][(y < x \lif \lnot \Atom{P}{y})])])].
\]
Deze formule zegt dat elke niet-lege verzameling een kleinste element heeft.
We behandelen een ``verzameling'' daarbij als een eenplaatsige relatie. Ook
de samenhang van een graaf kunnen we uitdrukken. Een samenhangende graaf kan
niet in twee niet-lege delen worden gesplitst zonder verbinding tussen die
delen. Dat staat in de volgende formule:
\[
\lnot \lexists[A][(\lexists[x][A(x)] \land \lexists[y][\lnot A(y)]
  \land \lforall[w][\lforall[z][((\Atom{A}{w} \land \lnot \Atom{A}{z})
      \lif \lnot \Atom{R}{w,z})]])].
\]
Als opgave kun je proberen de eindige !!{structure}s te definiëren waarvan
het !!{domain} een even aantal elementen heeft. Er kan nog meer: we kunnen de
reële getallen categorisch beschrijven als een volledig geordend lichaam dat
de rationale getallen bevat.

Tweede-ordelogica kan dus veel meer uitdrukken dan eerste-ordelogica. Dat is
het goede nieuws. Het slechte nieuws is dat geen effectief !!{derivation}
systeem volledig is voor de volledige tweede-ordesemantiek. Dat hadden we al
genoemd. Ook veel andere eigenschappen van de eerste-ordelogica gaan verloren,
waaronder compactheid en de stellingen van L\"owenheim--Skolem.

We kunnen de volledige tweede-ordesemantiek ook opgeven en de zwakkere
semantiek gebruiken. Dan is het minimale tweede-orde!!{derivation} systeem
wel volledig. Lees $\Proves$ als ``bewijsbaar in het minimale systeem'' en
$\Entails$ als ``volgt logisch in de zwakkere semantiek''. We kunnen dan
bewijzen dat uit $\Gamma \Entails !A$ steeds $\Gamma \Proves !A$ volgt. Willen we
bepaalde comprehensieaxioma's aan het !!{derivation} systeem toevoegen, dan
moeten we alleen structuren toelaten waarin die axioma's waar zijn. Laat
$\Delta$ bijvoorbeeld een verzameling comprehensieaxioma's zijn, misschien
zelfs de verzameling van alle comprehensieaxioma's. Dan geldt: als
$\Gamma \cup \Delta \Entails !A$, dan
$\Gamma \cup \Delta \Proves !A$. Is $!A$ met de gekozen
comprehensieaxioma's niet bewijsbaar, dan bestaat er dus een model van
$\lnot !A$ waarin die comprehensieaxioma's toch allemaal waar zijn.

Waarom geldt de volledigheidsstelling voor de zwakkere semantiek? Dat zien we
het gemakkelijkst als we tweede-ordelogica behandelen als logica met
verschillende soorten objecten, oftewel meersoortige logica. Eén soort krijgt
het gewone ``eerste-orde'' !!{domain}. Voor elke $k$ nemen we daarnaast
!!a{domain} met ``$k$-plaatsige relaties''. De taal krijgt ingebouwde
relatiesymbolen ``$\Atom{\Obj{true}_k}{R,x_1,\dots,x_k}$''. Dit symbool zegt
dat de relatie $R$ geldt voor $x_1$, \dots,~$x_k$. Hier is $R$ een variabele
van de soort ``$k$-plaatsige relatie''. De objecten $x_1$, \dots,~$x_k$
horen bij de eerste-ordesoort.

Na deze identificatie is de zwakke tweede-ordesemantiek in feite de gewone
semantiek voor meersoortige logica. We hebben al gezien dat meersoortige
logica in eerste-ordelogica kan worden ingebed. Als we de vertalingen heen en
terug buiten beschouwing laten, blijkt de zwakkere tweede-ordelogica dus
eerste-ordelogica in vermomming te zijn. Het !!{domain} bevat zowel
``objecten'' als ``relaties'', met axioma's die regelen hoe die twee soorten
zich gedragen.

\end{document}
